\section{Finite-module realizations and local-field geometry}
\label{sec:setup}

The first analytic goal is to construct a finite regular realization
for each choice of an omitted left action on the tensor copies of the
correspondence. Each realization includes the other left actions and
the diagonal right action. These realizations supply the finite traces
used to construct boundary measurements in Section~\ref{sec:links}.
We first remove the compact part of one conjugation action, then form
the tensor copies and the induced transports for the omitted source
action. Conjugation by these transports gives a genuine action of its
ambient group on the finite algebra. The remaining subsections describe
the geometric data to be read from the source-lattice Fourier indices.

Fix a nonzero
bifinite $L_\mu(\Gamma)$--$L_\omega(\Lambda)$ correspondence $K$ as in
Theorem~\ref{thm:arithmetic-exhaustion}. A right action
of $L_\omega(\Lambda)$ is written as the projective representation
\[
 \rho(h)\xi=\xi v_h^*,\qquad
 \rho(h)\rho(k)=\overline{\omega(h,k)}\rho(hk).
\]
The left representation is denoted by $\pi(g)\xi=u_g\xi$.  Thus $\pi$
and $\rho$ commute, although their multipliers need not be trivial.
Conjugation by either representation is an ordinary group action.
All Hilbert spaces and von Neumann algebras considered below are
separable in the senses appropriate to their normal representations.


\subsection{Compact reduction and finite-index return}

We will prove the geometric reconstruction first for an actual ICC
lattice $\Gamma<G$ and under the condition
\begin{equation}
 \Ad\rho\ \text{on}\
 L^2\bigl(\pi(L_\mu(\Gamma))',\operatorname{tr}\bigr)
       \ominus\mathbb C1
 \quad\text{has no nonzero finite-dimensional subrepresentation}.
 \label{eq:setup-weak-mixing}
\end{equation}
Here the commutant is a finite factor, equipped with its normalized
trace.  The following reductions justify this choice.  Their final use,
after the geometric partition has been constructed, is made in
Section~\ref{sec:recovery}.

\begin{lemma}[Twisted regular absorption]
\label{lem:setup-twisted-absorption}
Let $\alpha,\nu$ be normalized scalar cocycles on a countable group
$L$, and let $v$ be a $\nu$-projective unitary representation on a
Hilbert space $V$.  Then
\[
 \lambda_\alpha\otimes v
       \simeq \lambda_{\alpha\nu}\otimes1_V.
\]
If $L$ is ICC and $V$ is finite-dimensional, tensoring a finite normal
$L_\alpha(L)$-module with $v$ gives a finite normal
$L_{\alpha\nu}(L)$-module.  Any additional operators commuting with
every $v_h$ are preserved by the absorption identification.
\end{lemma}

\begin{proof}
On $\ell^2(L)\otimes V$ use the unitary
\[
 T(\delta_g\otimes\xi)=\delta_g\otimes v_g^*\xi.
\]
Indeed,
\begin{align*}
 T(\lambda_\alpha(h)\otimes v_h)T^*
          (\delta_g\otimes\xi)
 &=\alpha(h,g)\delta_{hg}\otimes v_{hg}^*v_hv_g\xi\\
 &=\alpha(h,g)\nu(h,g)\delta_{hg}\otimes\xi.
\end{align*}
A finite normal module over a finite factor embeds in finitely many
copies of the regular module
\cite[Theorems~10.1.1 and~10.2.1]{Jones2009}.  When $L$ is ICC both twisted group
algebras are factors.  Applying this unitary to such an embedding proves the second
assertion.  The formula also proves the last assertion directly.
\end{proof}

\begin{lemma}[Compact commutant reduction]
\label{lem:setup-compact-reduction}
Let $K$ be a nonzero bifinite
$L_\mu(\Gamma)$--$L_\omega(\Lambda)$ correspondence, with both groups
ICC.  There is a finite-index subgroup $\Lambda_0\leq\Lambda$ such
that the restriction of $K$ to the right subgroup factor is a finite
orthogonal sum of blocks
\[
 K_c=V_c\otimes K'_c,
 \qquad
 \pi(g)|_{K_c}=1\otimes\pi'_c(g),
 \qquad
 \rho(h)|_{K_c}=v^c_h\otimes\rho'_c(h)
       \quad(h\in\Lambda_0).
\]
The spaces $V_c$ are finite-dimensional.  The $v^c_h$ are normalized
projective matrix representations, with multipliers $\nu_c$, and
\begin{equation}
 \rho'_c(h)\rho'_c(k)
  =\frac{\overline{\omega(h,k)}}{\nu_c(h,k)}\rho'_c(hk).
 \label{eq:setup-reduced-right-multiplier}
\end{equation}
Each $K'_c$ is a finite normal module on both sides and satisfies
Equation~\eqref{eq:setup-weak-mixing} for its reduced representations.
There is no assertion that any $v^c$ has finite image.
\end{lemma}

\begin{proof}
Put $M=L_\mu(\Gamma)$, $N=L_\omega(\Lambda)$, and
$P=\pi(M)'$.  Finite $M$-module dimension makes $P$ a finite factor.
Moreover $N^{\mathrm{op}}\subset P$ has finite index.  To see the
latter without assuming the endpoint actions are mutual commutants,
regard $K$ as the equivalence module for the mutually commuting factors
$\pi(M)$ and $P$.  It has finite nonzero $P$-dimension.  Finitely many
copies of it therefore contain $L^2(P)$ as a $P$-module.  Their
restrictions to $N^{\mathrm{op}}$ are finite, by the original right
finiteness of $K$.  Consequently $L^2(P)$ has finite
$N^{\mathrm{op}}$-dimension, which is the index assertion.  We use
the standard finite-module and index facts in this argument in the
normal, trace-preserving convention; see
\cite{Jones1983} and \cite[Proposition~2.1]{PimsnerPopa1986}.

Let $D_c$ be the bounded elements of $P$ whose $\Ad\rho$-orbits are
precompact in $L^2(P)$.  The compact vectors of a unitary representation
form its closed almost-periodic subspace: they are the closed span of
its finite-dimensional subrepresentations.  Bounded compact elements
are closed under adjoints and products, since
\[
 \|xy-x'y'\|_2
 \leq \|x\|_\infty\|y-y'\|_2
       +\|x-x'\|_2\|y'\|_\infty
\]
on uniformly bounded sets.  Their bounded balls are strongly closed,
because bounded strong convergence in a finite von Neumann algebra
implies $L^2$ convergence.  They therefore form a von Neumann
subalgebra.  They are dense in the compact Hilbert subspace: the
nearest-point projection in $L^2(P)$ onto each closed operator-norm
ball is equivariant and nonexpansive, so maps compact orbits to compact
orbits, and these projections approximate every $L^2$ vector.

Twisted conjugation on $L^2(N)\ominus\mathbb C1$ has no nonzero
finite-dimensional subrepresentation.  Indeed, the diagonal of a
finite-rank invariant orthogonal projection in the Fourier basis is
constant on conjugacy classes: scalar conjugation phases disappear
from the diagonal.  The diagonal is summable.  Every nonidentity
conjugacy class is infinite, so all its entries vanish.  Only the
identity vector remains.  The trace-preserving expectation
$E_{N^{\mathrm{op}}}:P\to N^{\mathrm{op}}$ is equivariant and
$L^2$-contractive.  It follows that
\[
 E_{N^{\mathrm{op}}}(x)=\operatorname{tr}(x)1
             \quad(x\in D_c).
\]
By the finite-index expectation inequality there is $\varepsilon>0$
such that $E_{N^{\mathrm{op}}}(x)\geq\varepsilon x$ for $x\geq0$.
Thus a nonzero projection $q\in D_c$ has
$\operatorname{tr}(q)\geq\varepsilon$.  A finite algebra with such a
lower bound has only finitely many orthogonal nonzero projections,
and is finite-dimensional.

Restrict $\Lambda$ to the finite-index subgroup $\Lambda_0$ fixing
all the central blocks of $D_c$.  A central block $z$ now reduces both
endpoint actions.  Write
\[
 zK=V_c\otimes K'_c,
 \qquad zD_c=\mathcal B(V_c)\otimes1.
\]
The left action commutes with $D_c$, hence is $1\otimes\pi'_c$.
Choose normalized matrix implementers $v^c_h$ for
$\Ad\rho(h)|_{zD_c}$.  Uniqueness of an implementer up to scalar
gives
$v^c_hv^c_k=\nu_c(h,k)v^c_{hk}$, with $\nu_c$ a normalized cocycle.
After removing these implementers, $\rho(h)$ belongs to the
commutant of $\mathcal B(V_c)$ and therefore has the asserted tensor
form.  Multiplication gives
Equation~\eqref{eq:setup-reduced-right-multiplier}.

Normality and finiteness of $\pi'_c$ follow by taking a matrix corner
of the finite normal left module.  For the right action, first embed
the original restricted $\rho$ in a finite multiple of the
$\overline\omega|_{\Lambda_0}$-regular representation.  This is
possible because finite-index restriction preserves finite module
dimension.  Tensor with $\overline{v^c}$.  Lemma~
\ref{lem:setup-twisted-absorption} gives a finite normal module with
multiplier $\overline{\nu_c}\,\overline\omega$.  In
$\overline V_c\otimes V_c$ the vector corresponding to the identity
matrix is fixed by $\overline{v^c_h}\otimes v^c_h$.  Tensoring that
line with $K'_c$ exhibits $\rho'_c$ as a summand of the preceding
normal finite representation.  This proves its normality for exactly
the multiplier in Equation~\eqref{eq:setup-reduced-right-multiplier}.
Equivalently, its right factor has cocycle
$\omega|_{\Lambda_0}\nu_c$.

Finally, a compact vector in the reduced commutant embeds as
$1_{V_c}\otimes\xi$ on $zK$ and as zero on the other blocks.  It is
compact for $\Lambda_0$, and hence for $\Lambda$: the latter orbit
is contained in a finite union of translates of the former compact
orbit closure.  It therefore belongs to $L^2(D_c)$.  Its intersection
with the reduced commutant on this block is scalar.  This proves
Equation~\eqref{eq:setup-weak-mixing}.
\end{proof}

\begin{lemma}[Finite-index restriction and return]
\label{lem:setup-finite-index-induction}
Let $\Gamma_0\leq\Gamma$ and $\Lambda_0\leq\Lambda$ have finite
index, and restrict both cocycles literally to these subgroups.  Every
$M$--$N$ correspondence is a summand of the two-sided induction of
its restriction.  The inclusions are given, on the respective sides,
by
\begin{align}
 J_L\xi&=[\Gamma:\Gamma_0]^{-1/2}
       \sum_{r\Gamma_0\in\Gamma/\Gamma_0}
              u_r\otimes\pi(u_r)^*\xi,
       \label{eq:setup-left-induction}\\
 J_R\xi&=[\Lambda:\Lambda_0]^{-1/2}
       \sum_{\Lambda_0s\in\Lambda_0\backslash\Lambda}
              (\xi v_s^*)\otimes v_s.
       \label{eq:setup-right-induction}
\end{align}
These are bimodular isometries, with all scalar cocycles included.
Restriction and induction preserve bifiniteness.  An ICC group stays
ICC on finite-index restriction.
\end{lemma}

\begin{proof}
The $u_r$ in Equation~\eqref{eq:setup-left-induction} are an
orthonormal right basis for $L^2(M)$ over the subgroup factor, because
\[
 E_{L_{\mu|\Gamma_0}(\Gamma_0)}(u_r^*u_{r'})
   =\delta_{r,r'}1.
\]
The analogous left-basis identity for the $v_s$ is
$E_{L_{\omega|\Lambda_0}(\Lambda_0)}(v_sv_{s'}^*)
=\delta_{s,s'}1$.  The relative tensor-product inner products give
the asserted isometries.

For explicit left equivariance, if $gr=r'a$ with $a\in\Gamma_0$,
then
\[
 u_gu_r=\frac{\mu(g,r)}{\mu(r',a)}u_{r'}u_a.
\]
Move $u_a$ across the balanced tensor product.  The identity
$u_{r'}^*u_g=(\mu(g,r)/\mu(r',a))u_au_r^*$ shows that the term is
exactly the $r'$-term of $J_L(\pi(u_g)\xi)$.  The right calculation
is the same calculation in the opposite algebra.  Explicitly, when
$sh=at$ with $a\in\Lambda_0$, one has
\[
 v_sv_h=\frac{\omega(s,h)}{\omega(a,t)}v_av_t,
 \qquad
 v_hv_t^*=\frac{\omega(s,h)}{\omega(a,t)}v_s^*v_a.
\]
Balancing $v_a$ proves right equivariance of
Equation~\eqref{eq:setup-right-induction}.  The other endpoint action
commutes with every operation.  Thus the two maps compose to a
bimodular isometry into two-sided induction, whose range is a closed
summand.  The finite basis formulas also show preservation of finite
module dimensions.

If an element of $\Gamma_0$ has finite conjugacy class in
$\Gamma_0$, its conjugacy class in $\Gamma$ is a finite union of
translates of that class.  ICC of $\Gamma$ forces the element to be
the identity.  This proves the last assertion.
\end{proof}

For a source in $\mathscr K$, choose the actual isomorphism of
finite-index subgroups in Definition~\ref{def:class}.  Its image is
a finite-index subgroup of a reference lattice, hence itself a lattice
in the same ambient group.  Pull the scalar cocycle through that
isomorphism, without modifying it by a finite kernel.  The restrictions
retain property~(T) and ICC, as well as bifiniteness; property~(T)
permanence is used in its usual finite-covolume and finite-index forms
\cite[Theorem~1.7.1]{BHV2008}; see also the original finite-covolume
statement in \cite[Theorem~3]{DelarocheKirillov1968}.  Lemmas~\ref{lem:setup-compact-reduction} and
\ref{lem:setup-finite-index-induction} now reduce the geometric work
to the actual-lattice, weakly mixing setting.  Restoring $V_c$ after a
partition of $K'_c$ tensors every atom with $V_c$ and restores the
original exact right multiplier.  It changes neither the atom
permutations nor their freeness.  This observation will be used only
after that partition is proved.


\subsection{Tensor copies and complementary regular realizations}

From now until Section~\ref{sec:recovery}, let $\Gamma<G$ be an
actual ICC lattice and impose Equation~\eqref{eq:setup-weak-mixing}.
For a comparison of two source copies we use two copies of $K$; for
two comparisons sharing their first copy we use three. The right
comparison group acts diagonally. We omit one left source action from
the regular model so that its induced transports remain available on
$G/\Gamma$.


Take $m=1$ or $2$, and put
\[
 H=K^{\otimes(m+1)},\qquad
 R=\Gamma_1\times\cdots\times\Gamma_m\times\Lambda.
\]
The slots of $H$ are numbered $0,\ldots,m$; $\pi_j$ denotes the
left representation in slot $j$.  Let $W$ be the representation of
$R$ given by the $\pi_j$, $j\geq1$, and diagonal $\rho$.  Its split
multiplier is
\begin{equation}
 \alpha_R\bigl((g_1,\ldots,g_m,h),(g'_1,\ldots,g'_m,h')\bigr)
   =\prod_{j=1}^m\mu(g_j,g'_j)
        \overline{\omega(h,h')}^{\,m+1}.
 \label{eq:setup-tensor-multiplier}
\end{equation}

\begin{lemma}[Complementary regularity]
\label{lem:setup-normal-tensors}
The representation $W$ is normal and finite over the twisted factor
for $R$.  The same assertion holds with any one of the slots omitted
instead of slot zero.  The joint commutant of all $\pi_j(\Gamma)$,
$0\leq j\leq m$, and diagonal $\rho(\Lambda)$ is scalar.
\end{lemma}

\begin{proof}
Embed the omitted slot in a finite regular multiple for its
$\overline\omega$-projective right representation.  At that regular
index absorb the diagonal right actions of the remaining slots by
Lemma~\ref{lem:setup-twisted-absorption}.  This use of absorption
allows an infinite-dimensional auxiliary tensor space.  The
absorbing unitaries are formed from the other slots' $\rho(h)$ and
therefore commute with all their $\pi_j(g)$.  After absorption,
the $\Lambda$ action is regular with multiplier
$\overline\omega^{\,m+1}$, and each remaining $\pi_j$ can
independently be embedded in a finite regular multiple for
$L_\mu(\Gamma_j)$.  The resulting whole representation is a
subrepresentation of a finite regular multiple for exactly
Equation~\eqref{eq:setup-tensor-multiplier}.  The construction is
symmetric in the omitted slot.

Write $P=\pi(M)'$.  The final joint commutant is
$(P^{\overline\otimes(m+1)})^{\Ad\rho^{\otimes(m+1)}}$.
Condition~\eqref{eq:setup-weak-mixing} makes this fixed algebra
scalar.  For completeness, an invariant tensor with a factor in
$L^2(P)\ominus\mathbb C1$ gives a Hilbert--Schmidt intertwiner
between that factor's representation and another unitary
representation.  The positive compact operator obtained by composing
with its adjoint has a nonzero finite-dimensional invariant spectral
space, contrary to weak mixing.  Decompose each tensor factor as
$\mathbb C1\oplus(L^2(P)\ominus\mathbb C1)$ to apply this
argument to every nonconstant summand.
\end{proof}

Choose a column realization
\[
 H=eL^2(B)^k,\qquad e\in M_k(B),\qquad
 d=(\Tr_k\otimes\tau_B)(e)>0.
\]
Here $B$ is the opposite of the $\alpha_R$-twisted factor of $R$.
Using the adjoints of the opposite canonical generators identifies
it with the $\overline{\alpha_R}$-twisted group factor: write these
generators as $w_r$.  Thus $W_r$ is right multiplication by $w_r^*$,
which indeed has multiplier $\alpha_R$.

Set
\begin{equation}
 Y=G/\Gamma,\qquad
 \mathcal H=L^2(Y;H),\qquad
 \mathcal N=L^\infty(Y)\overline\otimes eM_k(B)e.
 \label{eq:setup-induced-model}
\end{equation}
Give $Y$ its invariant probability measure.  The trace on
$\mathcal N$ is
\[
 \operatorname{tr}_{\mathcal N}(A)
    =d^{-1}\int_Y(\Tr_k\otimes\tau_B)(A(y))\,dy,
\]
whereas the Hilbert norm of a column uses the unnormalized entrywise
trace sum.  All finite-matrix estimates permit rectangular products.
Whenever a family of coefficient algebras is varied, coefficient
traces are normalized and matrix sizes are fixed, or the corresponding
Hilbert-norm factors are retained explicitly.  An operator-norm bound
alone is not a bound independent of arbitrary trace rescaling.

Choose a Borel section $d_0:Y\to G$ and write
\[
 c(a,y)=d_0(ay)^{-1}ad_0(y),\qquad
 X_a(y)=\pi_0(c(a,y)),\qquad
 (U_af)(ay)=X_a(y)f(y).
\]
These transports compose up to a scalar function of modulus one on
the base.  More explicitly,
\[
 X_a(by)X_b(y)
   =\mu\bigl(c(a,by),c(b,y)\bigr)X_{ab}(y).
\]
Conjugation on base-decomposable operators is therefore a genuine
action, denoted by $\sigma=\Ad U$.

\begin{proposition}[Induced finite algebras]
\label{prop:setup-induced-algebras}
The transports $U_a$ are strongly continuous in $a$, and $\sigma$
preserves the trace on $\mathcal N$.  For $1\leq j\leq m$, let
$R'_j$ be $R$ with $\Gamma_j$ omitted, put
$\beta_j=\Ad W_{\Gamma_j}$, and let $\mathcal A_j$ be the
commutant of the base multipliers and $W|_{R'_j}$.  Then
\[
 (\mathcal A_j)^{\beta_j}=\mathcal N,
 \qquad
 (\mathcal A_j)^\sigma
 \ \text{has a faithful finite $\beta_j$-invariant trace},
 \qquad
 \mathcal N^G=\mathbb C1.
\]
\end{proposition}

\begin{proof}
Translation is continuous in measure on measurable maps on the
finite-measure quotient.  Applied to the Borel section, this shows
that $c(a,y)$ varies continuously in measure as a discrete-valued
map.  One can verify the assertion first on compact sets where the
section is continuous, by Lusin approximation, and then remove sets
of arbitrarily small measure.  The corresponding projective
unitaries are bounded, so their multiplication operators, together
with the base translations, vary strongly on $\mathcal H$.
Conjugation of the corner is inner at each fibre and translation
preserves base probability.  Thus $\sigma$ preserves the displayed
trace.  The first fixed-algebra identity is simply the commutant of
all of $W$ in the decomposable algebra.

For the second identity, undo induction.  If $A(y)$ is a
$\sigma$-invariant decomposable field, lift it, for
$g=d_0(y)\gamma$, as
\[
 \widetilde A(g)=\pi_0(\gamma)^*A(y)\pi_0(\gamma).
\]
The cocycle covariance makes this field left $G$-invariant.  It is
therefore almost everywhere constant, with value $T$, and right
$\Gamma$-covariance says that $T$ commutes with $\pi_0(\Gamma)$.
It also commutes with $W|_{R'_j}$.  Conversely every such $T$ gives
a fixed field.  Consequently $(\mathcal A_j)^\sigma$ is the
commutant of the complementary regular representation obtained by
omitting slot $j$.  Lemma~\ref{lem:setup-normal-tensors} makes this
commutant finite; its normalized trace is $\beta_j$-invariant,
since the remaining slot acts by unitaries in that commutant.
Including all the slots and using the last assertion of that Lemma
gives $\mathcal N^G=\mathbb C1$.
\end{proof}


\subsection{The source groups and full versus partial mixing}

The induced algebra has only scalar whole-group invariants. To turn
that fact into limits along Cartan paths, we need the mixing theorem
for the actual root-generated groups. We record the group conventions
here; the length and matrix coordinates are introduced afterwards.
We use the standard parabolic and root-group structure,
Cartan--Iwasawa decompositions, central-cover facts, root-generated
simplicity, and highest-weight theory
\cite{Borel1991,BT1965,BruhatTits1972,BruhatTits1984,Tits1964,Humphreys1972}.

\begin{convention}[Factors and panel groups]
\label{conv:setup-panels}
At a complex place the local group is regarded as a real Lie group.
The noncompact property-(T) archimedean factors are the higher-rank
ones and the groups locally isomorphic to $Sp(n,1)$, $n\geq2$, or
$F_{4(-20)}$.  At a finite place the factors have relative rank at
least two. For higher-rank property~(T), see
\cite[Theorem~6]{DelarocheKirillov1968}; the real rank-one
classification and its Kostant attribution are recalled in
\cite[Section~3.3 and Theorem~3.5.4]{BHV2008}. Relative rank one
at a finite place gives an unbounded action on the Bruhat--Tits tree
and is incompatible with property~(T) \cite[Remark~1.6.3]{BHV2008}.

In a rank-one semisimple Levi position exactly one almost simple
factor is isotropic.  Its minimal boundary is the full panel residue.
The compact and toric normal factors of the Levi act trivially on
this boundary.  Source transports use its root-generated group.
When an algebraic central cover is needed, use the subgroup generated
by its lifted unipotent radicals.  The root-group isomorphisms below
map this subgroup onto the effective root-generated group with finite
central kernel.  Every application of Howe--Moore takes invariants
for this root-generated group or lift.
Projectivities use the effective adjoint group; the adjoint point
group is allowed for changes of frame.  We never replace a full flag
by an oriented or topological covering flag.
\end{convention}

To check these conventions, take two opposite parabolics.  Their
refinement spaces identify with the flags of the common Levi.  On
the transverse open cell, the radical is removed by projection, and
Levi transport acts by its inner algebraic automorphisms.  The
anisotropic normal factors are compact and have no flag coordinate
in the relative building; the central torus acts trivially there as
well.  Simply connected covers have finite central kernel, contain
the root-generated groups in their images on points, and have image
of at most finite index in the adjoint point group.  They have the
same Bruhat flag orbits.  At real places any remaining ineffective
or component extension is compact or finite.

For the cover assertion, a central covering induces an isomorphism
on the corresponding relative root groups, so its image contains
the subgroup generated by the unipotent radicals
\cite[Theorem~2.20(iii) and Proposition~2.24(ii)]{BT1972Complements}.
Its kernel on local points is finite.  For a central kernel $C$,
the rational-point quotient injects into $H^1(k,C)$, which is finite
over the present local fields
\cite[Propositions~1.12 and~1.17, Theorem~6.1]{BorelSerre1964}.
This proves the finite-index assertion directly; no stronger
identification of all rational points with the root-generated
group is needed.


\begin{lemma}[Product mixing and partial escape]
\label{lem:setup-product-mixing}
For the action on $L^2(\mathcal N)$, escape in every simple factor
gives weak operator convergence to the projection onto constants.
If a path is the identity outside a subproduct $G_0$ and escapes in
every factor of $G_0$, its weak limit is the projection onto
$L^2(\mathcal N)^{G_0}$, equivalently the conditional expectation
onto $\mathcal N^{G_0}$.  On the orthogonal complement of a simple
factor's invariants, its matrix-coefficient decay remains uniform
after composition with arbitrary commuting group unitaries.
\end{lemma}

\begin{proof}
Use the Howe--Moore theorem for the noncompact isotropic almost
simple root-generated groups and their root-generated lifts from
Convention~\ref{conv:setup-panels} \cite{HoweMoore1979}; an independent proof with this root-generated group convention is given in \cite[Theorem~1.1 and Definition~4.23]{Ciobotaru2015}.
We first justify the type-I
input for these actual groups.  For an isotropic simple factor
$\mathbf H$ over $k$, choose a nonzero parabolic unipotent radical
$\mathbf U$.  The $k$-span of the $\mathbf H(k)$-conjugates of
$\operatorname{Lie}(\mathbf U)$ is a nonzero invariant Lie ideal,
hence all of $\operatorname{Lie}(\mathbf H)$.  Choose a basis
$X_1,\ldots,X_d$ from these conjugate subspaces.  The map
\[
 (t_1,\ldots,t_d)\longmapsto\prod_{j=1}^d\exp(t_jX_j)
\]
has invertible differential at zero and takes values in
$\mathbf H(k)^+$.  The local-field inverse function theorem makes
this subgroup open; at finite places see
\cite[Proposition~1.11]{Gloeckner2008}.  The same proof applies to
the root-generated subgroup of an algebraic simply connected cover.
At an archimedean place, the generating unipotent groups are connected,
so the root-generated group is connected.  Being open as well, it is
the identity component of the ambient real algebraic group; for a
complex group, use restriction of scalars to $\mathbb R$.  The same
applies to its algebraic lift.  These groups are therefore type~I by
\cite[Theorem~1]{Dixmier1957}.  At a finite place, the ambient reductive
group of local points is type~I \cite{Bernstein1974}, and type~I passes
to its open root-generated subgroup
\cite[Proposition~6.E.21(1)]{BekkaHarpe2019}.

Restrict a separable unitary representation $\pi$ to one simple
factor $S$.  Its type-I decomposition with multiplicity is
\[
 \mathcal H=\int_{\widehat S}^{\oplus}
       (H_\xi\otimes M_\xi)\,d\eta(\xi),\qquad
 \pi(s)=\int_{\widehat S}^{\oplus}(\xi(s)\otimes1)\,d\eta(\xi).
\]
In this central decomposition, every unitary commuting with $S$
acts fibrewise as $1\otimes V_\xi$
\cite[Propositions~5.4.7 and~8.6.4, Theorem~8.6.6]{Dixmier1977};
see also \cite[Theorem~7.32]{Folland2015}.  This is the decomposition
of the restriction to $S$, so no invariance of a finer decomposition
of the other factors is required.  For finite tensor sums in a
nontrivial fibre, uniformly over the unitaries $V_\xi$,
\[
 \left|\left\langle(\xi(s)\otimes V_\xi)
       \sum_j e_j\otimes u_j,\sum_k f_k\otimes v_k\right\rangle\right|
 \leq\sum_{j,k}|\langle\xi(s)e_j,f_k\rangle|
                    \|u_j\|\,\|v_k\|.
\]
Howe--Moore makes the right side tend to zero along escape.
Approximation by finite tensor sums is uniform because all operators
have norm one.  For arbitrary fixed fibre vectors the bound
$\|a_\xi\|\,\|b_\xi\|$ is integrable.  Dominated convergence
therefore proves the stated uniform decay off the $S$-invariants.

Let $P_i$ project onto the invariants of the $i$th simple factor.
These projections commute with every group action and with each
other.  Expand the identity as the finite sum of the mutually
orthogonal products obtained from $\prod_i(P_i+(1-P_i))$.
In every sector containing $1-P_i$, the uniform decay just proved
kills its coefficients when that factor escapes, even while all
other factor translations vary.  Thus for all-factor escape the
weak limit is $\prod_iP_i$, which is the projection onto constants
by Proposition~\ref{prop:setup-induced-algebras}.  For partial escape,
expand only over the factors of $G_0$; the other translations are the
identity.  The limit is the projection onto
$L^2(\mathcal N)^{G_0}$, namely the trace-preserving conditional
expectation.  The same argument applies to the simple Levi factors
used later.
\end{proof}


\subsection{Panel groups, Cartan coordinates, and Haar volumes}

The following geometric facts will specify the Fourier-index data used
by the boundary measurements. At a panel, only one noncompact rank-one
Levi factor acts on the residue; on a Cartan path, exact root coordinates
will determine which flag coordinates survive.

An archimedean panel inside a higher-rank component cannot have
noncompact Lie algebra $\mathfrak f_{4(-20)}$: its complexification
would have to occur as a component of a proper absolute Dynkin
subdiagram, and the Dynkin list has no such occurrence.  At a finite
place the effective rank-one Lie algebra is simple also over
$\mathbb Q_p$.  Indeed after restriction of scalars its absolutely
simple factors form one Galois orbit, so a $\mathbb Q_p$-ideal is
either zero or the whole algebra.  An adjoint automorphism centralizing
an open subgroup centralizes all sufficiently small inner
exponentials, and hence this Lie algebra; it is therefore the
identity.  These assertions concern the ambient local groups, not a
discreteness assumption on a lattice projection.

Let $\Pi$ be the union of the simple \emph{indivisible} relative
roots of the factors.  A face of type $I\subset\Pi$ has strictly
positive simple coordinates exactly in $I$; opposition of types
is denoted by a star.  Choose a special maximal compact in each
factor.  Cartan and Iwasawa decompositions can first be made on the
simply connected covers and then descended.  Write $\mu(g)$ for the
Cartan vector only when its argument is a local-group element; this
use of $\mu$ is distinct from the scalar two-cocycle
$\mu(g,h)$.  For a positive root $\alpha$, put
\[
 m_\alpha=\dim_{k}\mathfrak g_\alpha
     \quad\hbox{at a finite place},
 \qquad
 2\rho=\sum_{\alpha>0}m_\alpha\alpha,
\]
and use the analogous real restricted-root multiplicities at
archimedean places.  Double roots are included in this sum.  The
fundamental coweight rays $e_i$ are dual to the simple indivisible
roots.

\begin{convention}[Exact finite-place translation coordinates]
\label{conv:setup-exact-lengths}
At a finite place use the module absolute value, with
$|\varpi|=q^{-1}$.  The Cartan vector is the actual dominant
translation vector of a Cartan representative in the split
centralizer: its pairing with every centralizer character is the
exact logarithm of that character's module.  We do not replace this
by a coarsely equivalent length in identities involving heights.
All displacements and offsets along $e_i$ range on a sufficiently
divisible mesh in this logarithmic translation lattice.  Cocharacter
rationality and common multiples of relative coroot translations
and central isogenies ensure that these elements and their
semisimple projections to every standard Levi have lifts in the
corresponding root-generated covering groups.  There are only finitely
many standard Levis, so one divisible mesh suffices.  At real places
the corresponding split exponentials lie in the generated identity
lift.  No argument at a finite place requires every real
parameter to be an actual split displacement.
\end{convention}

Bounded parts of the centralizer act equicontinuously, with bounded
inverse, on the relevant weight spaces.  This permits compact frames
while keeping the character normalization exact.

\begin{lemma}[Haar size of Cartan shells]
\label{lem:setup-haar-shells}
For a bounded-thickness shell about a dominant vector $L$, its Haar
measure is at most $C\exp(2\rho(L))$.  For one fixed sufficiently
large thickness, there is also a lower bound
$c\exp(2\rho(L))$.  The constants are independent of $L$ on the
permissible Cartan mesh.  The analogous assertions hold for the
finite product, with the sum of the factors' $2\rho$.
\end{lemma}

\begin{proof}
At a real place use the Cartan density, a product of the appropriate
powers of $\sinh\alpha(L)$.  The upper bound follows on every fixed
bounded neighborhood.  For the lower bound, even when $L$ lies on a
wall, choose inside the shell one fixed bounded offset box whose
simple coordinates are bounded away from zero.  On this box every
positive-root factor is bounded below by a fixed multiple of its
exponential.  The box has fixed positive Euclidean volume.

At a finite place the Cartan double coset has measure
$[K:K\cap aKa^{-1}]$ times the fixed compact measure.  This index is
comparable to $\exp(2\rho(L))$.  One can see the uniform comparison
in compact root charts: a sufficiently small analytic compact open
subgroup has product coordinates in negative-root, centralizer, and
positive-root directions.  Under conjugation by the dominant
translation, its intersection is obtained by shrinking the expanding
root lattices.  The index is the product of their module indices,
with exponent $\sum_{\alpha>0}m_\alpha\alpha(L)$.  Equivariant
exponentials on small root lattices and parabolic product uniqueness
give this calculation; changing the compact frame changes the small
subgroups by a uniformly bounded index.  Other Cartan representatives
are within bounded distance of the chosen split mesh.  This proves
both bounds there.  Multiplication of Haar measures gives the product
assertion.
\end{proof}

Fixed left or right multiplication changes vector distances by a
bounded amount.  We use ordinary compact-rotation probabilities on
flags, or equivalently the Haar measure classes in their root
coordinates.  Proper algebraic incidence degeneracies have measure
zero.  All incidence and Weyl positions are those of the relative
spherical building.  A prescribed reduced longest gallery from a
chamber parametrizes its open opposition cell by successive panel
coordinates, with both measure-null directions preserved; this is
the ordered unipotent factorization of the Bruhat cell.  The same
statement holds at finite places.


\subsection{Compactified matrices and admissible scalar tests}

The measured geometric index will always lie in a copy $\Gamma_j$ of
the source lattice, inside its known ambient local group. We now give
a compact matrix description of the flags and residual projectivities
that can survive as that index escapes. This description is algebraic
at finite places; only its scalar observables will act on our complex
Hilbert spaces.


For each relative vertex type choose a highest-line representation
as follows.  Take the determinant line of the relevant parabolic
nilradical in an exterior power and the cyclic submodule it generates.
Over a splitting field, choose a compatible positive system.  This
line is a highest line, with a weight strictly positive on the omitted
absolute nodes and zero on the complementary coroots.  Its relative
weight is a positive multiple of the relevant fundamental weight.
The cyclic highest-weight module is absolutely irreducible and its
top split space is one-dimensional.  Lowering along an omitted
absolute simple root produces a weight separated by a positive
multiple of just the corresponding relative simple root.  Every
other weight drops by a nonnegative combination involving that root.
Consequently a normalized Cartan matrix has rank-one limit exactly
when that simple gap diverges.

At a finite place these are $k$-rational modules, used only to form
compact projective matrix data.  Use compact-invariant lattice norms
and compact sup-norm frames, recording both the matrix and its inverse
up to projective normalization.  At an archimedean target, for Hilbert
highest-line amplifications use a real proximal module irreducible
on complexification, with a compact-invariant inner product.  Taking
higher tensor weights is allowed.  No Hilbert space is tensored with
a vector space over a nonarchimedean field.

\begin{lemma}[Face blocks and residual projectivities]
\label{lem:setup-face-blocks}
Let a sequence have diverging simple gaps of types $D$, and keep its
complementary gaps in a bounded window.  Its compactified matrix data
determine an image $D$-flag, a reverse $D^*$-flag, and an invertible
residual Levi transport between their transverse refinement spaces.
These data are independent of compact frame choices.  On a deeper
stratum the already determined image vertices persist.  On compact
transverse windows the residual rank-one panel lengths differ by
bounded amounts from the actual omitted Cartan gaps.
\end{lemma}

\begin{proof}
Pass to convergent compact Cartan frames.  In each highest-weight
module, precisely the weights with maximal value on the open face
survive projective normalization.  They are the submodule generated
from the highest line by the complementary Levi.  Bounded
complementary gaps keep the surviving block and its inverse bounded
in those frames.  Thus the surviving block acts by an invertible Levi
transformation.  The same analysis of the inverse matrix identifies
the reverse face.

Equivalently, write an ordinary flag in the open Bruhat cell
transverse to the reverse face.  Conjugating its radical coordinates
by the Cartan translation contracts the coordinates belonging to the
diverging gaps; the complementary Levi coordinates remain.  First
project to the reverse refinement space, then apply the surviving
Levi transformation.  This is the intrinsic residual projectivity.
A change of compact frame conjugates these coordinates by the
corresponding compact Levi projectivity and leaves the intrinsic map
unchanged.  Passing to a deeper stratum preserves all preceding image
vertices.  In a face omitting one type in a higher-rank component,
the surviving rank-one split displacement is exactly that omitted
coordinate; bounded frames change its length by a uniformly bounded
amount.  The argument works simultaneously for one omitted type in
several components.

In an archimedean module the surviving top block contains the highest
lines of all incident flags, and those lines span it by the Levi
action.  At a finite place this is an algebraic statement about the
$k$-module only; the later operator-algebra argument instead uses
scalar clopen ranges.  Exhausting buffered transverse patches gives
the assertions on the full stratum.
\end{proof}

Absolute growth on a highest line is measured with the compact-equivariant
norm.  Relative to the exact Cartan exponent, its correction is the
test on the input flag direction.  On a transverse cell the normalized
ratio converges to the corresponding absolute pairing.  At a finite
place the character valuations of the nonsplit centralizer give the
same statement with Convention~\ref{conv:setup-exact-lengths}; an
arbitrary bounded-error replacement of the character exponent is not
made.

\begin{convention}[Smooth and clopen symbols]
\label{conv:setup-symbols}
On real coordinate windows, a smooth symbol has the specified
uniformly bounded finite number of derivatives, increased whenever
a separated expansion is needed.  On finite-place coordinate boxes,
its scalar-test counterpart is locally constant at a uniformly fixed
clopen level.  The underlying coordinate maps themselves remain
analytic.  Uniformity means that normalized denominators are bounded
away from zero on a buffered larger window, so that all the chosen
clopen levels pull back to a common level.  In the bounded-width
case of Equation~\eqref{eq:input-2}, a test of one argument may
instead be an arbitrary bounded function of that argument.
\end{convention}

Here is the elementary expansion fact implicit in this convention.
On a fixed product of buffered real coordinate boxes, extend the
symbol smoothly to a product of tori.  Fourier coefficients, after
more differentiations than the total coordinate dimension, are
bounded by a fixed summable sequence.  On a finite-place box there
are only finitely many cells at the prescribed level.  Combining
Fourier characters with their cell indicators expresses a joint
symbol as an absolutely summable sum of products of unary symbols.
For a uniformly bounded family of smooth seminorms the coefficient
bound is the stronger uniform bound
\[
       \sum_n\sup_\theta |c_n(\theta)|<\infty,
\]
where $\theta$ denotes any additional parameter.  This is the bound
needed for the Gram-sequence estimates later.

On a buffered transverse matrix-product window, projective
normalization has a uniformly nonzero denominator.  Matrix
multiplication and its inverse-coordinate version therefore have
uniform smooth bounds at real places and uniform clopen levels at
finite places.  At a finite place the logarithm of a weight-matrix
product is the sum of its two entry logarithms and a bounded locally
constant correction of uniform level.  Additional logarithmic phases
are kept as unary factors on the two entries; only the correction is
expanded in the compact coordinates.  Growing radial cutoffs are
used in the unshifted regimes specified in
Section~\ref{sec:heights}.  In particular this convention makes no
claim that an arbitrarily translated radial cutoff confines a
rescaled logarithm to a fixed compact box.

Common panel equalities will be used as exact substitutions by finite
clopen partitions in the finite-place slots, followed by regular
smooth division in the real slots.  Countably many determining
window and angular identities are established before multiplying by
arbitrary Borel output projections.  One need not assume a uniform
transverse estimate for arbitrary discontinuous joint masks.

\begin{lemma}[Interpolation in an affine building]
\label{lem:setup-building-interpolation}
Fix an endpoint $x$ in one of the affine buildings above.  In an
apartment containing $x,y$, multiply the dominant simple coordinates
of $y-x$ by independent numbers $t_i\in[0,1]$.  The resulting point
is independent of the apartment and depends on $y$ with a uniform
Lipschitz constant.  The same uniform assertion holds for variation
of $x$, after opposition and replacement of the fractions by their
complements.  Constants for a finite product depend only on its root
systems.
\end{lemma}

\begin{proof}
Write $L_i=\alpha_i(y-x)\geq0$ and let $T_t$ denote the linear map
with simple coordinates $t_iL_i$.  For every positive root
$\alpha=\sum_i n_i\alpha_i$, where $n_i\geq0$,
\[
 0\leq\alpha(T_t(y-x))
      =\sum_i n_it_iL_i
      \leq\sum_i n_iL_i=\alpha(y-x).
\]
Thus the new point belongs to every affine root half-apartment
containing both endpoints: it lies in their root enclosure.
Apartment intersections are enclosed.  The apartment transition
maps agree on this enclosure, and the Weyl stabilizer of a wall
displacement also fixes its transformed displacement, since its zero
simple coordinates stay zero.  This proves independence.

On any apartment and Weyl sector, $T_t$ has a uniformly bounded
Euclidean operator norm, independent of the fractions.  Subdivide
a segment of varying endpoints into cells and Weyl sectors relative
to $x$, and choose apartments containing $x$ and those cells.
Local finiteness makes this a finite subdivision on a compact
segment.  Sum the Euclidean bounds and use apartment independence
at the interfaces.  Reversing the endpoints applies opposition and
the complementary fractions, proving the other Lipschitz bound.
\end{proof}

We use successive limits along sequences, realized by free
ultrafilters when convenient.  A faster time is always sent to
infinity with all slower parameters fixed.  Contracted-radical
estimates are invoked only in that order.  Noncompact isotropic Levi
factors contain discrete nonabelian free subgroups: choose two
rank-one translations with separated attracting and repelling data
and sufficiently large powers.  The usual ping-pong domains, in the
rank-one symmetric space or the locally finite tree, give a discrete
Schottky subgroup.  The full chamber action is a topologically
amenable homogeneous-space action because its minimal parabolic is
amenable; the anisotropic part is compact modulo the central torus.
Restriction to a closed subgroup preserves this amenability
\cite[Propositions~2.2 and~2.5, Example~2.7(5)]{AnantharamanDelaroche2002}.  These facts are used for the
source local group and its lattice, not for an action of the unknown
comparison group on a building.


\begin{convention}[Operators and Fourier coordinates]
\label{conv:setup-operators}
An invariance under $U$ is always a conjugation invariance of
decomposable operators, so base scalar phases do not change it.
An ``other group'' in a commutant includes $\Lambda$ unless
$\Lambda$ is explicitly the measured group.  A target Fourier arrow
in the geometric arguments records $w_{(h,\cdot)}$ and the Cartan
or flag data of $h\in\Gamma$ in a measured source-lattice copy;
it does not record geometric data of the $\Lambda$ index.  A bounded
column is an operator-bounded tracial matrix column, essentially
uniformly bounded when it depends on $y$.  The closed range of a
positive effect means its support space when a range inclusion is
asserted.
\end{convention}
